\chapter{绪论}
\label{chap:introduction}

\section{研究背景与问题提出}

知识图谱问答将自然语言问题与结构化知识连接起来，使答案能够由明确的查询操作获得。对于包含条件筛选、集合关系和数值比较的问题，模型不仅需要识别问题中的实体与属性，还需要保持操作之间的依赖。本文关注约三十亿参数的小模型：在模型规模和推理预算受到约束时，如何通过可执行查询及其反馈，提高问答过程的可靠性，并将这一方法用于具有来源与条件要求的雷达资料。

执行器能够返回中间结果、空集合或失败信息，为模型提供继续求解的依据。然而，提供反馈与学会利用反馈是两个不同的问题。若训练只包含顺利到达答案的参考轨迹，模型在自由生成时进入偏离参考的状态后，可能缺少与该状态相适应的后续行为。单纯增加普通成功轨迹，也未必能够覆盖这种训练与推理之间的差异。因此，本文从真实运行状态出发，研究怎样构造可以继续执行的恢复目标，并在控制监督预算的条件下检验其作用。

雷达领域使这一问题具有明确的应用约束。同一份资料可能同时描述整机、天线与显示部件；相同属性可能随工作模式和测量条件取不同数值；若干离散选项也不能直接合并成连续区间。问题中的型号能够检索到，并不意味着其参数已被正确归属；查询成功也不保证回答保留了必要限定。此外，回答还应能够追溯到实际给定的资料，不能以合法但不支持断言的引用代替证据。由此，领域应用既需要可执行的问答方法，也需要保留条件与来源的数据表示，以及区分不同失败环节的评价。

\section{研究目标与研究问题}

本文的目标是在有限模型规模下，研究真实执行状态与后续监督之间的联系，并以公开知识图谱任务和雷达来源问答分别检验方法效果及应用边界。围绕这一目标，提出三个研究问题。

第一，如何将雷达原文整理成可查询、可定位的来源读法，同时避免将结构一致误认为事实已核验？这一问题关注主体、部件、值类型、适用条件与来源之间的对应关系，也关注问题、参考和推理证据的分离。整理结果必须能回到原始资料复查，并明确区分机器检查、AI 审阅和真人核验的证据等级。

第二，在相同初始模型、训练问题和监督 token 预算下，引入真实偏离状态及其可执行恢复后缀，能否提高小模型知识图谱问答的正确率与执行完成率？这一问题是本文的主要训练研究。对照需要使普通续训和恢复续训共享问题及目标后缀，避免将额外训练量当成恢复状态的作用；同时，首次偏离参考程序只用于定位样本，不能直接当作语义错误标签。

第三，公开任务上的增量能否延续到雷达应用；若不能，领域链路中的哪些限制值得单独研究？为回答这一问题，本文分别检查接口使用、记录选择、支持页覆盖、正文正确和引用支持。对于检索后仍然答错的现象，再以固定语义记录比较平铺与条件绑定表示，从而考察证据组织的作用，而不预设它属于恢复训练的迁移收益。

\section{研究内容与技术路线}

本文围绕上述问题组织三项相互衔接的工作。第一项是来源约束的雷达数据整理与评价材料构建。以固定的制造商资料为依据，保存完整页块、结构化读法及引用定位，对多主体共享信息保留共同归属，对缺失或有歧义的内容保留状态。八来源主体材料共含 502 条读法与 96 道中文问题，每组十二题；读法及问题均经过 AI 编写或审阅，记录数不等于独立事实数，相关来源内的问题也不视为独立来源样本。

第二项是执行状态与恢复后缀的配对监督方法。先由初始分步模型在训练题上自由运行，从最终失败的轨迹中选择具有共同成功前缀的样本。普通组从此前缀执行参考后缀；恢复组额外保留首次分歧动作及真实观察，再重映射中间结果句柄，使同一参考后缀适用于当前状态。构造后逐步重放验证，训练时只监督后缀中的模型输出，不要求模型模仿分歧动作或生成工具观察。两组混入相同的普通轨迹，并匹配逐记录监督 token 数。该方法的研究改进在于状态选择、后缀构造与监督位置控制，所用语言模型、参数高效微调及交叉熵优化均为已有技术。

第三项是分层评价与雷达应用验证。公开实验以普通配对续训为主要对照，以一次性程序生成提供质量与成本参照。领域实验先用共同合成材料适配有限查询接口，再依据新来源评价实际回答；随后另设八来源对照，在不训练新模型的条件下比较原文、平铺记录和条件绑定记录。平铺与绑定解码后得到完全相同的声明语义记录，保留顺序、重复项与字段类型。此约束不意味着结构化记录包含原文的全部信息，也不消除整理成本和输入长度差异。

\begin{table}[htbp]
\centering\small
\caption{本文技术路线及主要验证对象}
\label{tab:intro-route}
\begin{tabular}{p{0.20\linewidth}p{0.36\linewidth}p{0.32\linewidth}}
\toprule
研究环节 & 主要处理 & 验证对象 \\
\midrule
来源与数据 & 保留主体、条件、量值和出处，分离问题与参考 & 来源定位、表示完整性及审阅等级 \\
执行反馈训练 & 真实分歧状态、句柄重映射与配对后缀监督 & 匹配监督预算下的问答增量 \\
应用与表示评价 & 共同接口适配、固定证据及分项审阅 & 迁移边界、正文与引用的不同作用 \\
\bottomrule
\end{tabular}
\end{table}

三项工作共同形成从资料到执行、再到答案与证据的研究链路，其中只有恢复配对监督承担主要训练方法贡献。来源表示提供可检查的数据基础，应用评价给出方法的适用证据；二者不另行包装为两项新的学习算法。

\section{证据范围与论文组织}

本文在 KQA Pro 官方训练数据中构建项目训练、开发和留出划分。最终 2,000 题项目留出评价中，恢复组相对普通组的两个续训种子增量为 4.65 和 4.85 个百分点，预定平均增量为 4.75 个百分点。两个种子共享初始模型和冻结语料，属于续训随机性的配对检查，并非两次全流程独立训练；项目留出集也不是官方隐藏测试集。相同监督预算不代表相同计算量，分步方法与一次性程序生成的成本差异需同时报告。

雷达实验保留未获得收益的结果：共同适配后的四来源评价没有建立恢复组的额外回答优势。八来源表示实验每路均保留 96 题分母，平铺与绑定的正文正确数为 73 和 80、题级绑定错误为 17 和 11。联合正确要求正文正确、所给证据支持正文且引用支持断言；原始计数为 17 和 52，原文路线为 56。统一精确引用解析后，平铺与绑定变为 55 和 68，原文仍为 56；该分析在原结果之后登记，不改正文，也不产生新的模型预测。因此，联合指标的全部改善不能归于事实推理。上述评价仍依赖 AI 参考与审阅；另备的 24 题真人核验材料截至本稿完成数为零，不能提升为人工金标。

全文共七章。第一章提出研究问题、内容与证据范围。第二章介绍知识图谱问答、执行反馈学习及领域证据组织的相关基础。第三章说明雷达数据整理、来源表示与评价材料构建。第四章给出恢复状态与执行后缀配对监督方法。第五章报告公开任务上的受控实验、执行成本及机制诊断。第六章分析雷达接口适配、来源问答和固定证据表示对照。第七章归纳研究结论、限制与后续方向，使方法效果、领域应用及尚未完成的核验各有对应依据。
